% Einheit 2 von 4 – Wahrscheinlichkeit einstufig (bank/wahrscheinlichkeit/e2.jsonl, zusammenbau v0.1)
%% TODO Zweigzeile Teil 1: was hier gelernt wird (Satz in Schülersprache aus dem Einheitstitel) – Einheit 2
\einheitenkopf[e2]{Einheit 2 von 4 · Wahrscheinlichkeit einstufig}
\zweigzeile{ab Kl. 6, je nach Buch bis Kl. 7 · P10 oft}
\setcounter{aufgabe}{10}

%% TODO Titel: Ich-kann-Satz fehlt in der Bank (Platzhalter: Kettenname) – Nr. 11
%% TODO Anweisung: ein Satz über der ersten Teilaufgabe fehlt in der Bank – Nr. 11
\begin{aufgabe}{Wie viele Stufen?}
\begin{teile}
\teil Tim dreht ein Glücksrad einmal. Wie viele Stufen hat der Versuch? Kreuze an, rechne nichts. \\ \kreuz{eine Stufe} \\ \kreuz{zwei Stufen} \\ \kreuz{drei Stufen}
\end{teile}
\end{aufgabe}

%% TODO Titel: Ich-kann-Satz fehlt in der Bank (Platzhalter: Kettenname) – Nr. 12
%% TODO Anweisung: ein Satz über der ersten Teilaufgabe fehlt in der Bank – Nr. 12
\begin{aufgabe}{Gleich groß?}
\begin{teile}
\teil Das Glücksrad hat fünf gleich große Felder. Welcher Anteil der Felder ist rot? Kreuze an. \\ \kreuz{$\frac{1}{3}$} \\ \kreuz{$\frac{3}{5}$}
\end{teile}
\end{aufgabe}

%% TODO Titel: Ich-kann-Satz fehlt in der Bank (Platzhalter: Kettenname) – Nr. 13
%% TODO Anweisung: ein Satz über der ersten Teilaufgabe fehlt in der Bank – Nr. 13
\begin{aufgabe}{Einstufig}
\begin{teile}
\teil In einer Lostrommel liegen 30 Nieten und 5 Gewinne. Gesucht ist die Wahrscheinlichkeit für einen Gewinn. Unterstreiche, was möglich und was günstig ist; rechne nicht. möglich: \leerfeld , günstig: \leerfeld
\end{teile}
%% TODO Darstellung für schwach fehlt in der Bank (wahrscheinlichkeit-e2-k3-s1-v1; Abschnitt „Für schwache Schüler“)
\swz{Ein Würfel wird einmal geworfen. Wie groß ist die Wahrscheinlichkeit für eine 4?}{}
%% TODO Darstellung für schwach fehlt in der Bank (wahrscheinlichkeit-e2-k3-s1-v2; Abschnitt „Für schwache Schüler“)
\swz{Ein Würfel mit acht gleichen Flächen trägt die Zahlen 1 bis 8. Wie groß ist die Wahrscheinlichkeit für eine 7?}{}
%% TODO Darstellung für schwach fehlt in der Bank (wahrscheinlichkeit-e2-k3-s1-v3; Abschnitt „Für schwache Schüler“)
\swz{Ein Würfel mit vier gleichen Flächen trägt die Zahlen 1 bis 4. Wie groß ist die Wahrscheinlichkeit für eine 2?}{}
%% TODO Darstellung für schwach fehlt in der Bank (wahrscheinlichkeit-e2-k3-s1-v4; Abschnitt „Für schwache Schüler“)
\swz{Ein Würfel trägt auf seinen sechs Flächen die Buchstaben A, B, C, D, E und F. Wie groß ist die Wahrscheinlichkeit für E?}{}
%% TODO Darstellung für schwach fehlt in der Bank (wahrscheinlichkeit-e2-k3-s1-v5; Abschnitt „Für schwache Schüler“)
\swz{Ein Würfel hat sechs verschiedenfarbige Flächen, eine davon ist grün. Wie groß ist die Wahrscheinlichkeit für Grün?}{}
\swfrage{Was bleibt gleich, was ändert sich?}
\swz{Das Glücksrad hat acht gleich große Felder. Wie groß ist die Wahrscheinlichkeit für Blau?}{\kreisdiagramm{rot/1,blau/1,rot/1,gelb/1,rot/1,blau/1,gelb/1,rot/1}}
%% TODO Darstellung für schwach fehlt in der Bank (wahrscheinlichkeit-e2-k3-s3-v1; Abschnitt „Für schwache Schüler“)
\swz{In einer Urne liegen 5 rote, 3 blaue und 1 grüne Kugel. Eine Kugel wird gezogen. Wie groß ist die Wahrscheinlichkeit für Blau?}{}
%% TODO Darstellung für schwach fehlt in der Bank (wahrscheinlichkeit-e2-k3-s4-v1; Abschnitt „Für schwache Schüler“)
\swz{Ein Würfel wird einmal geworfen. Wie groß ist die Wahrscheinlichkeit für eine Zahl kleiner als 3?}{}
%% TODO Darstellung für schwach fehlt in der Bank (wahrscheinlichkeit-e2-k3-s5-v1; Abschnitt „Für schwache Schüler“)
\swz{In einer Schachtel liegen 50 Murmeln: 18 blaue, 9 rote und 23 grüne. Eine Murmel wird gezogen. Wie groß ist die Wahrscheinlichkeit für Rot als Bruch, als Dezimalzahl und in Prozent? \hfill (P10 2019 OS)}{}
%% TODO Darstellung für schwach fehlt in der Bank (wahrscheinlichkeit-e2-k3-s6-v1; Abschnitt „Für schwache Schüler“)
\swz{Bei einer Tombola liegen 45 Nieten und 15 Gewinnlose in der Trommel. Wie groß ist die Wahrscheinlichkeit für einen Gewinn? \hfill (P10 2014 OS)}{}
%% TODO Darstellung für schwach fehlt in der Bank (wahrscheinlichkeit-e2-k3-s7-v1; Abschnitt „Für schwache Schüler“)
\swz{Ein Glücksrad hat acht gleich große Felder mit den Zahlen 1 bis 8. Wie groß ist die Wahrscheinlichkeit für weder 1 noch 8? \hfill (P10 2014 OS)}{}
%% TODO Darstellung für schwach fehlt in der Bank (wahrscheinlichkeit-e2-k3-s8-v1; Abschnitt „Für schwache Schüler“)
\swz{Ein Glücksrad hat 4 gleich große Felder in Rot und Gelb. Die Wahrscheinlichkeit für Gelb ist 75\,\%. Wie viele Felder sind gelb? \hfill (P10 2018 OS)}{}
\swa{Zeichne ein Glücksrad mit 6 gleich großen Feldern, bei dem die Wahrscheinlichkeit für Rot $\frac{1}{2}$ und für Blau $\frac{1}{3}$ ist; die übrigen Felder sind gelb.}{\kreisleer[2]}
%% TODO Darstellung für schwach fehlt in der Bank (wahrscheinlichkeit-e2-k3-s10-v1; Abschnitt „Für schwache Schüler“)
\swz[3]{Auf einem Teller liegen 18 Muffins, 3 davon mit Chili, die anderen mit Schokolade. Jonas hat schon 2 Schoko-Muffins gegessen. Mara sagt: „Die Wahrscheinlichkeit, jetzt einen Chili-Muffin zu erwischen, ist $\frac{3}{18}$.“ Hat sie Recht? Begründe. \hfill (P10 2018 OS)}{}
%% TODO Darstellung für schwach fehlt in der Bank (wahrscheinlichkeit-e2-k3-s11-v1; Abschnitt „Für schwache Schüler“)
\swz{Ein Würfel wird 600-mal geworfen, 93-mal fällt eine 6. Vergleiche die relative Häufigkeit mit der Wahrscheinlichkeit.}{}
%% TODO Darstellung für schwach fehlt in der Bank (wahrscheinlichkeit-e2-k3-s12-v1; Abschnitt „Für schwache Schüler“)
\swz[3]{Nils sagt: „Beim Würfeln ist eine 6 seltener als eine 1.“ Stimmt das? Begründe.}{}
\swz[3]{Bei einem Schulfest gibt es Lose mit den Nummern 201 bis 700. Einen Trostpreis gibt es für jedes Los, dessen Nummer auf 45 endet, außer für das Los 445 (Hauptgewinn). Timo hat ein Los und sieht nur die letzte Ziffer: 5. Er sagt: „Die Wahrscheinlichkeit für einen Trostpreis ist $\frac{4}{50}$.“ Hat er Recht? Begründe. \hfill (P10 2016 OS)}{}
\end{aufgabe}

%% TODO Titel: Ich-kann-Satz fehlt in der Bank (Platzhalter: Kettenname) – Nr. 14
%% TODO Anweisung: ein Satz über der ersten Teilaufgabe fehlt in der Bank – Nr. 14
\begin{aufgabe}{Topf mit P = 50 \% auswählen}
\begin{teile}
\teil Drei Dosen enthalten je 8 Bonbons: links 5 Kirsch und 3 Zitrone, in der Mitte 4 Kirsch und 4 Zitrone, rechts 2 Kirsch und 6 Zitrone. Aus welcher Dose zieht man Kirsch mit der Wahrscheinlichkeit 50\,\%? Kreuze an. \hfill (P10 2015 OS) \\ \kreuz{linke Dose} \\ \kreuz{mittlere Dose} \\ \kreuz{rechte Dose}
\end{teile}
\end{aufgabe}

%% TODO Titel: Ich-kann-Satz fehlt in der Bank (Platzhalter: Kettenname) – Nr. 15
%% TODO Anweisung: ein Satz über der ersten Teilaufgabe fehlt in der Bank – Nr. 15
\begin{aufgabe}{erwartete Anzahl bei n Versuchen}
%% TODO Darstellung für schwach fehlt in der Bank (wahrscheinlichkeit-e2-k5-s1-v1; Abschnitt „Für schwache Schüler“)
\swz{Ein Würfel wird 300-mal geworfen. Wie oft erwartest du ungefähr eine 5?}{}
\end{aufgabe}

%% TODO Titel: Ich-kann-Satz fehlt in der Bank (Platzhalter: Kettenname) – Nr. 16
%% TODO Anweisung: ein Satz über der ersten Teilaufgabe fehlt in der Bank – Nr. 16
\begin{aufgabe}{Einstufig – Fehler finden, Begründen, Darstellungswechsel, Anwendung}
\swz[3]{In einer Packung sind 50 Batterien, 3 davon sind leer. Paul gibt die Wahrscheinlichkeit, eine leere zu ziehen, mit $\frac{3}{47}$ an. Finde den Fehler.}{}
\swfrage{Beim Glücksrad nimmt das rote Feld den halben Kreis ein, die beiden anderen Felder sind kleiner. Tom sagt: „Die Wahrscheinlichkeit für Rot ist $\frac{1}{3}$, weil es drei Felder gibt.“ Begründe, warum das falsch ist.}
\swz{Die Gewinnwahrscheinlichkeit beträgt 35\,\%. Schreibe sie als gekürzten Bruch und als Dezimalzahl.}{}
\swz[3]{Ein Hersteller wirbt: „In jedem 8. Überraschungsei steckt eine Sammelfigur.“ Jonas kauft 40 Eier. Wie viele Figuren kann er erwarten? Kann er sich darauf verlassen?}{}
\end{aufgabe}

%% TODO Merkkasten Einheit 2: 6 Zeilen, Vorgabe höchstens fünf (3.1)
\uebersichtskasten{Wahrscheinlichkeit: günstige Ergebnisse geteilt durch alle möglichen Ergebnisse – wenn alle Ergebnisse gleich wahrscheinlich sind (Laplace). \\ Glücksrad mit 12 gleich großen Feldern, 5 rot, 4 blau, 3 gelb: P(rot) = 5/12 ≈ 0,42 ≈ 42 \% \\ Ereignis aus mehreren Ergebnissen: die günstigen Felder zusammenzählen. \quad P(blau oder gelb) = (4 + 3)/12 = 7/12 \\ Gegenereignis: P(nicht A) = 1 − P(A). \quad P(nicht rot) = 1 − 5/12 = 7/12 – dasselbe wie „blau oder gelb“ \\ Rückwärts: Anzahl = Wahrscheinlichkeit mal Gesamtzahl. \quad P(grün) = 25 \% bei 16 Feldern → 4 grüne Felder \\ Relative Häufigkeit: Treffer geteilt durch Versuche; bei vielen Versuchen liegt sie nahe an der Wahrscheinlichkeit (Gesetz der großen Zahlen) – einzelne Versuche sagt sie nicht voraus.}
